\section{Euler tour}

\begin{frame}
\centering
\huge Euler Tour
\end{frame}

\begin{frame}
  \frametitle{Consultas de xor en caminos (sin updates)}

  \begin{itemize}
    \item Te dan un arbol con numeros en las aristas
    \item Querés soportar consultas de xor en caminos
    \item \(N \leq 2 \cdot 10^5\)
    \item \(Q \leq 2 \cdot 10^5\)
  \end{itemize}
  Esto se puede resolver descomponiendo en caminos pesados.

  \vspace{0.75em}

  Pero también hay una solución muchísimo más simple.
\end{frame}

\begin{frame}
  \frametitle{Consultas de xor en caminos: idea}

  La idea es construir el \textbf{Euler tour}: un array de los valores de aristas en el orden en que te los encontrás si hacés un DFS del árbol.

  \vspace{0.75em}

  \begin{itemize}
    \item Se inserta el valor tanto al \textbf{ir} como al \textbf{volver} de cada arista.
    \item Para responder una consulta de xor en un camino \(u\)--\(v\), consultamos en ese array. El rango corresponde a las posiciones de \(u\) y \(v\) en el recorrido.
  \end{itemize}
\end{frame}

\begin{frame}
  \frametitle{Consultas de xor en caminos: por qué funciona}

  Si mirás el intervalo de valores entre una aparición de \(u\) y una de \(v\), cualquier arista por la que ``fuimos y volvimos'' entre medio en el DFS se cancela, porque el xor es su propia inversa.

  \vspace{0.75em}

  Y coincidentemente esas aristas son las que \textbf{no} están en el camino \(u\)--\(v\). O sea, las que no se cancelan son exactamente las del camino \(u\)--\(v\).

  \vspace{0.75em}

  Por lo tanto, la respuesta es simplemente el xor de los valores de las aristas en el camino \(u\)--\(v\).
\end{frame}

\begin{frame}[fragile]
  \frametitle{Consultas de xor en caminos: código}

\begin{lstlisting}
struct edge {
    int u, v, x;
    int opp(int w) { return u + v - w; }
}
edge edges[maxn];
vector<int> g[maxn];
vector<int> euler_tour;
int position[maxn];
void dfs(int u, int p) {
    for (int i : g[u]) {
        edge e = edges[i];
        int v = e.opp(u);
        if (v == p) continue;

        position[u] = sz(euler_tour);
        euler_tour.push_back(e.x);
        position[v] = sz(euler_tour);
        dfs(v, u);
        euler_tour.push_back(e.x);
    }
}

void init() {
    dfs(0, -1);
    init_array(euler_tour);
}

void query(int u, int v) {
    int pu = position[u], pv = position[v];
    return xor_array(pu, pv);
}
\end{lstlisting}
\end{frame}
